\documentclass{exam}

\newif\ifA
\newif\ifkey
\newif\ifprintselected   % required by exam_config's \ansbox

%%%%%%%%%%%%%%%%%%%%%%%%%%%%%%%%%%%%%%%%%%%%%%%%%%%%%%
%%% SET THESE IF VARIABLES
%%%%%%%%%%%%%%%%%%%%%%%%%%%%%%%%%%%%%%%%%%%%%%%%%%%%%%
% Lab name -- labs are named by the lesson they belong to, never by a
% sequential lab number (Neil, 2026-10-12): this is the L8 Lab.
\def\labLesson{8}

% is this the key or not?
\ifdefined\iskey
	\ifnum\iskey=1
		\keytrue
	\else
		\keyfalse
	\fi
\else
	%% Pick one manually
	\keyfalse % if this is NOT the key
%	\keytrue     % if this IS the key
\fi

% set up the header and footer - change for every term
\ifdefined\thisTermIn
	\def\thisterm{\thisTermIn}
\else
	\def\thisterm{Spring 2027}
\fi
%
\def\thisclass{ECE 448}

%%%%%%%%%%%%%%%%%%%%%%%%%%%%%%%%%%%%%%%%%%%%%%%%%%%%%%
%%% END SET THESE IF VARIABLES
%%%%%%%%%%%%%%%%%%%%%%%%%%%%%%%%%%%%%%%%%%%%%%%%%%%%%%
% packages and shortcuts
\input{myPackages_gr}
\usepackage[hidelinks, linkcolor=red]{hyperref}
\input{myShortcuts}
\setlength{\parindent}{0pt}
\setlength{\parskip}{6pt}

%% exam class customization
\input{exam_config}
% ECE 448 document fonts.
%
% Body text is Barlow (SIL OFL; the faces and the license are in
% latex/fonts/Barlow/). Barlow is not in TeX Live, so it is loaded as an
% OpenType face through fontspec -- which means the document must be built
% with lualatex, not pdflatex. build_practice.sh and build_lab.sh both do.
%
% Compiling with pdflatex anyway is not an error: the guard below falls back
% to the LaTeX default so the source still compiles under any engine and can
% be syntax-checked. The resulting PDF just is not the one we ship.
%
% Paths are relative to latex/, which is where both build scripts run
% (each does cd "$(dirname "$0")"). Loading by path rather than by installed
% family name means no system font install and no fontconfig, so a fresh
% clone builds a byte-comparable PDF anywhere.
%
% MATH IS DELIBERATELY LEFT IN COMPUTER MODERN. Barlow has no math companion,
% and CM math under sans body text is already the house pairing: the decks set
% words in Source Sans Pro and let MathJax render symbols in its CM-derived
% face. Moving math to a sans as well would mean unicode-math plus a sans math
% font (Fira Math being the plausible one) and re-checking every equation in
% the course for symbol coverage. That is a separate decision, not a side
% effect of this one.

\usepackage{iftex}
\ifLuaTeX
	\usepackage{fontspec}
	\defaultfontfeatures{Ligatures=TeX}
	\setmainfont{Barlow}[
		Path           = fonts/Barlow/,
		Extension      = .ttf,
		UprightFont    = *-Regular,
		ItalicFont     = *-Italic,
		BoldFont       = *-Bold,
		BoldItalicFont = *-BoldItalic,
	]
	\setsansfont{Barlow}[
		Path           = fonts/Barlow/,
		Extension      = .ttf,
		UprightFont    = *-Regular,
		ItalicFont     = *-Italic,
		BoldFont       = *-SemiBold,
		BoldItalicFont = *-BoldItalic,
	]
	% microtype earns its place here specifically because of Barlow. Its
	% glyphs are narrower than Computer Modern's, so justified prose sets
	% tighter and TeX has to squeeze interword space to fit; two paragraphs
	% in the L02/L03 sets went overfull on the first Barlow build. Protrusion
	% and font expansion absorb that -- both are lualatex features, so this
	% belongs inside the guard with everything else.
	\usepackage{microtype}

	% Mono stays Latin Modern: Barlow ships no monospaced face, and verbatim
	% and code need one that actually is fixed-pitch.
\fi
   % Barlow body text (needs lualatex)
\printselectedfalse      % full worked solutions in the key, not answers-only
\CorrectChoiceEmphasis{\color{red}}
\SolutionEmphasis{\color{red}}
\renewcommand\questionlabel{\textbf{\thequestion.}}
\renewcommand{\thepartno}{\alph{partno}}
\renewcommand{\partlabel}{(\thepartno)}

% Fillable table cell: blank in the student copy, the expected value in red on
% the KEY. Fixed-width centered columns keep the blank cells writable.
\newcolumntype{K}[1]{>{\centering\arraybackslash}p{#1}}
\newcommand{\kv}[1]{\ifkey{\color{red}#1}\fi\rule[-0.10in]{0pt}{0.34in}}

% Answers are not in this public file. Every solution, \kv value, and \ansbox
% value is \keysol{<id>}, defined in the private ECE448-faculty repo at
% lab-keys/src/ECE448_Lab_L08_FSK_solutions.tex; build_lab.sh puts that
% folder on TEXINPUTS (KEYDIR) for the KEY. A KEY build without it stops.
\long\def\defkeysol#1#2{\expandafter\long\expandafter\def\csname keysol@#1\endcsname{#2}}
\def\keysol#1{\ifcsname keysol@#1\endcsname\csname keysol@#1\endcsname\fi}
\ifkey
	\IfFileExists{ECE448_Lab_L08_FSK_solutions.tex}%
		{\input{ECE448_Lab_L08_FSK_solutions}}%
		{\errmessage{KEY build needs ECE448_Lab_L08_FSK_solutions.tex from ECE448-faculty/lab-keys/src (set KEYDIR)}}
\fi

% print the answers if this is the key
\ifkey
	\printanswers % if this IS the key
\else
	\noprintanswers  %if this is NOT the key
\fi

% print the header and footer
% need to print key on top if this is the key
\ifkey
	\firstpageheader{}{\color{red}{ KEY}}{}
	\runningheader{\thisclass\hspace{1pt} -- \thisterm }{L\labLesson\ Lab\hspace{1pt}\color{red}{ KEY}}{\thepage\hspace{1pt} of\hspace{1pt} \numpages}
	\firstpagefooter{}{}{}
	\runningfooter{}{}{}
\else
	\firstpageheader{}{}{}
	\runningheader{\thisclass \hspace{1pt} -- \thisterm }{L\labLesson\ Lab\hspace{1pt}}{\thepage\hspace{1pt} of\hspace{1pt} \numpages}
	\firstpagefooter{}{}{}
	\runningfooter{}{}{}
\fi

\runningheadrule


\begin{document}
\pagestyle{headandfoot}   % exam class

%%%%%%%%%%%%%%%%%%%%%%%%%%%%%%%%%%%%%%%%%%%%%%%%%%%%%%
% title block
%%%%%%%%%%%%%%%%%%%%%%%%%%%%%%%%%%%%%%%%%%%%%%%%%%%%%%

\begin{center}
USAF ACADEMY \\
DEPARTMENT OF ELECTRICAL AND COMPUTER ENGINEERING \\
LESSON \labLesson\ LAB - FSK \\
\textbf{Lab Packet} \\
\thisterm \\[4mm]
\end{center}

\vspace{-4pt}
\textbf{Name:} \rule[-2pt]{2.6in}{0.4pt} \hfill
\textbf{Section:} \rule[-2pt]{0.9in}{0.4pt} \hfill
\textbf{Date:} \rule[-2pt]{1.1in}{0.4pt}

\vspace{4pt}

\fullwidth{\textbf{Learning Objective 4.1: } I can describe FSK modulation and implement
an FSK transmitter and receiver in GNU Radio.}

\vspace{2pt}

This packet is the L\labLesson\ Lab, in three parts, and you turn it in once, complete,
before Lesson 9. Do \textbf{Part 1} on your own before the lab period: run the link
flowgraph, read its display, and work out the block settings you will use at the
bench. \textbf{Part 2} is your own measurements, on your own laptop, recorded in
class. Do \textbf{Part 3} on your own after the lab: explain what you saw and answer
the analysis questions.

\textbf{Grading.} One submission, graded against Objective 4.1 on the course's
1 / 0 scale, with revision allowed under course policy. Paste screenshots where a
question asks for one.

\textbf{No hardware, no transmitting.} Every flowgraph in this lab runs on your laptop.
The Channel Model block stands in for the air, and nothing in this lab transmits.

The flowgraphs are on the L8 lesson page: \texttt{L08-fsk-tx.grc} and
\texttt{L08-fsk-link.grc}. Every number below assumes the flowgraphs' settings:
$f_s = 32$ kHz, 32 samples per symbol, and a line of 40 characters and a newline, 41 bytes.

\begin{questions}

\fullwidth{\large\textbf{Part 1: Pre-Lab} \normalsize(individual, before the lab period)}

%%%%%%%%%%%%%%%%%%%%%%%%%%%%%%%%%%%%%%%%%%%%%%%%%%%
\question \textbf{Run the link.} Open \texttt{L08-fsk-link.grc} in GRC and run it
with the sliders at their defaults.

\begin{parts}

	\begin{minipage}[t]{0.58\linewidth}
		\part After about 20 s, record the Lines label.
		\begin{solution}[0.5in]\keysol{sol1}\end{solution}
	\end{minipage}\hfill%
	\begin{minipage}[t]{0.37\linewidth}
		\raggedleft \vspace{12pt}
		\ansbox[1.6in]{\keysol{ans1}}
	\end{minipage}

	\part Paste a screenshot of the running flowgraph's window.
		\begin{solution}[1.6in]\keysol{sol2}\end{solution}

\end{parts}

%%%%%%%%%%%%%%%%%%%%%%%%%%%%%%%%%%%%%%%%%%%%%%%%%%%
\newpage
\question \textbf{Settings and readings.} Question 3 checks the bandwidth against your spectrum.

\begin{parts}

	\part Fill in the transmitter and receiver settings for each deviation.

	\begin{center}
	\renewcommand{\arraystretch}{1.25}
	\begin{tabular}{| l | K{0.95in} | K{0.95in} | K{0.95in} | K{0.95in} |}
	\hline
	$\Delta f$ & sensitivity $k$ & demod gain $g$ & bandwidth $B$ & LPF cutoff \\\hline
	1500 Hz & \kv{\keysol{kv1}} & \kv{\keysol{kv2}} & \kv{\keysol{kv3}} & \kv{\keysol{kv4}} \\\hline
	2000 Hz & \kv{\keysol{kv5}} & \kv{\keysol{kv6}} & \kv{\keysol{kv7}} & \kv{\keysol{kv8}} \\\hline
	\end{tabular}
	\end{center}

		\begin{solution}[0.9in]\keysol{sol3}\end{solution}

	\part Run \texttt{L08-fsk-link.grc}. On its spectrum, read how far the two tones
	stand above the noise floor with the noise slider at 0.5, then at 1.0. Describe
	what the eye diagram does between the two.

	\begin{center}
	\renewcommand{\arraystretch}{1.25}
	\begin{tabular}{| l | K{1.4in} | K{1.4in} |}
	\hline
	Noise voltage & $0.5$ & $1.0$ \\\hline
	Tones above floor (dB) & \kv{\keysol{kv9}} & \kv{\keysol{kv10}} \\\hline
	\end{tabular}
	\end{center}

		\begin{solution}[0.8in]\keysol{sol4}\end{solution}

	\begin{minipage}[t]{0.58\linewidth}
		\part Each line is a packet: a 64-bit access code, a 32-bit header, and the 41
		bytes of the line, its 40 characters and a newline. Find the bits per packet, the time per packet, and
		the number of packets in 30 s.
		\begin{solution}[1.0in]\keysol{sol5}\end{solution}
	\end{minipage}\hfill%
	\begin{minipage}[t]{0.37\linewidth}
		\raggedleft \vspace{12pt}
		bits = \ansbox[1.0in]{\keysol{ans2}}

		\vspace{4pt}
		time = \ansbox[1.0in]{\keysol{ans3}}

		\vspace{4pt}
		in 30 s = \ansbox[1.0in]{\keysol{ans4}}
	\end{minipage}


\end{parts}

%%%%%%%%%%%%%%%%%%%%%%%%%%%%%%%%%%%%%%%%%%%%%%%%%%%
\newpage
\fullwidth{\large\textbf{Part 2: Measurements} \normalsize(your own laptop, in class)}

\question \textbf{Build the transmitter.} Start GRC with an empty canvas. Do not open
\texttt{L08-fsk-tx.grc}: build it yourself as the lab page lists, from a Vector Source
of message bytes, Unpack K Bits, Chunks to Symbols, Repeat, Frequency Mod, a
Throttle, and a QT GUI Frequency Sink, with a QT GUI Range named \texttt{dev}.

\begin{parts}

	\part Paste a screenshot of your flowgraph on the GRC canvas. Every connection
	must be black, not red.
		\begin{solution}[1.6in]\keysol{sol6}\end{solution}

	\part Run it and record the frequencies of the two spectral peaks at each
	deviation. Paste a screenshot of the spectrum at 2000 Hz, and compare its width
	at each deviation with your bandwidth $B$ from Question 2(a).

	\begin{center}
	\renewcommand{\arraystretch}{1.25}
	\begin{tabular}{| l | K{1.4in} | K{1.4in} |}
	\hline
	Deviation (Hz) & $1500$ & $2000$ \\\hline
	Peaks (kHz) & \kv{\keysol{kv11}} & \kv{\keysol{kv12}} \\\hline
	\end{tabular}
	\end{center}

		\begin{solution}[1.6in]\keysol{sol7}\end{solution}

\end{parts}

%%%%%%%%%%%%%%%%%%%%%%%%%%%%%%%%%%%%%%%%%%%%%%%%%%%
\newpage
\question \textbf{Noise.} Run \texttt{L08-fsk-link.grc} with the frequency offset at 0.
For each noise voltage, set the slider, wait about 30 s, and record the Lines label.
The count restarts each time a slider moves. Sent is the time divided by your
packet time from Question 2(c).

\begin{parts}

	\part Fill in the table.

	\begin{center}
	\renewcommand{\arraystretch}{1.25}
	\begin{tabular}{| l | K{0.6in} | K{0.6in} | K{0.6in} | K{0.6in} | K{0.6in} | K{0.6in} |}
	\hline
	$\sigma$ & $0.5$ & $0.8$ & $1.0$ & $1.2$ & $1.5$ & $2.0$ \\\hline
	Time (s) & & & & & & \\\hline
	Sent & & & & & & \\\hline
	Received & \kv{\keysol{kv13}} & \kv{\keysol{kv14}} & \kv{\keysol{kv15}} & \kv{\keysol{kv16}} & \kv{\keysol{kv17}} & \kv{\keysol{kv18}} \\\hline
	Correct & \kv{\keysol{kv19}} & \kv{\keysol{kv20}} & \kv{\keysol{kv21}} & \kv{\keysol{kv22}} & \kv{\keysol{kv23}} & \kv{\keysol{kv24}} \\\hline
	\end{tabular}
	\end{center}

		\begin{solution}[0.6in]\keysol{sol8}\end{solution}

	\part Paste a screenshot of the eye diagram at $\sigma = 1.2$, and describe how it
	differs from the eye at $\sigma = 0.5$.
		\begin{solution}[1.6in]\keysol{sol9}\end{solution}

	\begin{minipage}[t]{0.58\linewidth}
		\part Set the noise back to 0.5 and raise it in steps of 0.1 until a wrong
		line appears. Record that noise voltage, and compare the tones' height above
		the floor there with your reading at 1.0 in Question 2(b).
		\begin{solution}[0.7in]\keysol{sol10}\end{solution}
	\end{minipage}\hfill%
	\begin{minipage}[t]{0.37\linewidth}
		\raggedleft \vspace{12pt}
		$\sigma$ = \ansbox[1.0in]{\keysol{ans5}}

		\vspace{4pt}
		vs.\ 2(b): \ansbox[1.0in]{\keysol{ans6}}
	\end{minipage}

\end{parts}

%%%%%%%%%%%%%%%%%%%%%%%%%%%%%%%%%%%%%%%%%%%%%%%%%%%
\newpage
\question \textbf{Frequency offset.} Set the noise voltage back to 0.5. For each
offset, set the slider, wait about 20 s, and record the two levels of the eye diagram
and the Lines label.

\begin{parts}

	\part Fill in the table.

	\begin{center}
	\renewcommand{\arraystretch}{1.25}
	\begin{tabular}{| l | K{0.75in} | K{0.75in} | K{0.75in} | K{0.75in} | K{0.75in} |}
	\hline
	Offset (Hz) & $0$ & $300$ & $600$ & $800$ & $1000$ \\\hline
	Levels & \kv{\keysol{kv25}} & \kv{\keysol{kv26}} & \kv{\keysol{kv27}} & \kv{\keysol{kv28}} & \kv{\keysol{kv29}} \\\hline
	Lines & \kv{\keysol{kv30}} & \kv{\keysol{kv31}} & \kv{\keysol{kv32}} & \kv{\keysol{kv33}} & \kv{\keysol{kv34}} \\\hline
	\end{tabular}
	\end{center}

		\begin{solution}[0.6in]\keysol{sol11}\end{solution}

	\part Fill the $\Delta f = 1000$ Hz row from Question 4(a) ($\sigma = 1.2$) and
	from one run at $\sigma = 0.5$ and an offset of 1200 Hz before changing
	\texttt{dev}. Then stop the flowgraph and change the \texttt{dev} variable to
	2000; the transmitter, the filter, and the demodulator all follow it. Run it at
	each setting below for about 20 s and record the Lines label.

	\begin{center}
	\renewcommand{\arraystretch}{1.25}
	\begin{tabular}{| l | K{1.5in} | K{1.5in} |}
	\hline
	Setting & $\sigma = 1.2$, offset 0 & $\sigma = 0.5$, offset 1200 Hz \\\hline
	$\Delta f = 1000$ Hz & \kv{\keysol{kv35}} & \kv{\keysol{kv36}} \\\hline
	$\Delta f = 2000$ Hz & \kv{\keysol{kv37}} & \kv{\keysol{kv38}} \\\hline
	\end{tabular}
	\end{center}

		\begin{solution}[0.6in]\keysol{sol12}\end{solution}

\end{parts}

%%%%%%%%%%%%%%%%%%%%%%%%%%%%%%%%%%%%%%%%%%%%%%%%%%%
\newpage
\fullwidth{\large\textbf{Part 3: Writeup} \normalsize(individual, after the lab)}

\question \textbf{Analysis.}

\begin{parts}

	\part From your table in Question 4 and your screenshots, describe in words what
	the eye and the received text did as you raised the noise from 0.5 to 2.0. Explain
	why a small turn of the slider takes the link from clean text to no text, and what
	that means for the link margin of Lesson 2.
		\begin{solution}[1.5in]\keysol{sol13}\end{solution}

	\part Take the largest offset at which your link still delivered correct lines in
	Question 5(a). Convert it to parts per
	million at 915 MHz, and explain what it says about running this link over the air
	between two SDRs whose crystals are accurate to a few ppm.
		\begin{solution}[1.3in]\keysol{sol14}\end{solution}

	\part Recommend a deviation of 1000 or 2000 Hz for this link, and justify it from
	your measurements in Questions 4 and 5. Name what the choice gives up.
		\begin{solution}[1.3in]\keysol{sol15}\end{solution}

\end{parts}

\end{questions}

\vspace{1.0em}
\noindent\textbf{Documentation:} \rule[-2pt]{0.6\linewidth}{0.4pt}

\end{document}
